% GENERATED FILE. Edit canonical lessons, then run npm run generate:books.
% canonical-source-hash: fnv1a64:283df847b38c2970
% canonical-lessons: PT-C82-poder, PT-C82-dever, PT-C82-necessidade, PT-C82-oficio, PT-C82-levar

\chapter{To be able to, To have to, Need, Trade, To take}
\label{ch:pt-to-be-able-to-to-have-to-need-trade-to-take}
\hlchaptermodality{82}

\begin{chapteropening}
\textbf{By the end of this chapter:} \emph{I can say to be able to, to have to, a need, a trade and to take.}

Complete the last lesson of chapter 82: I can say to be able to, to have to, a need, a trade and to take.
\end{chapteropening}

\section[poder]{poder — to be able to, can}

\label{lesson:PT-C82-poder}



\begin{quote}
Before the new one: say the Portuguese for to lose, then the Portuguese for to arrive.
\end{quote}

\subsection*{You'll want to know: poder}
\textbf{poder} — \textquotedblleft{}to be able to, can\textquotedblright{}.

\begin{cousinweb}
\textbf{poder} $\leftarrow$ Vulgar Latin \emph{*potere} (EtymDB 2.1), from Latin \emph{posse}, \textquotedblleft{}to be able\textquotedblright{} — inherited.
\end{cousinweb}

\begin{grammarlens}[title={What you've built}]
One more word for this chapter.
\end{grammarlens}

\subsection*{Your turn}
\begin{itemize}
\raggedright
  \item \emph{Say it:} \emph{poder}
  \item \emph{Say it:} \emph{poder}, once more
  \item \emph{Say it:} say \emph{poder}
  \item \emph{Recall:} say the Portuguese for to lose, then the Portuguese for to arrive, then say \emph{poder} again
\end{itemize}

\begin{tcolorbox}[breakable,colback=teal!4,colframe=teal!35!black,title={Before you move on}]
What does poder mean? (\textquotedblleft{}To be able to, can\textquotedblright{}.) Say it once more.
\end{tcolorbox}
\section[dever]{dever — to have to, must}

\label{lesson:PT-C82-dever}



\begin{quote}
Before the new one: say the Portuguese for to arrive, then the Portuguese for to be able to.
\end{quote}

\subsection*{You'll want to know: dever}
\textbf{dever} — \textquotedblleft{}to have to, must\textquotedblright{}.

\begin{cousinweb}
\textbf{dever} $\leftarrow$ Latin \emph{debere}, \textquotedblleft{}to owe\textquotedblright{} (EtymDB 2.1) — inherited.
\end{cousinweb}

\begin{grammarlens}[title={What you've built}]
One more word for this chapter.
\end{grammarlens}

\subsection*{Your turn}
\begin{itemize}
\raggedright
  \item \emph{Say it:} \emph{dever}
  \item \emph{Say it:} \emph{dever}, once more
  \item \emph{Say it:} say \emph{dever}
  \item \emph{Recall:} say the Portuguese for to arrive, then the Portuguese for to be able to, then say \emph{dever} again
\end{itemize}

\begin{tcolorbox}[breakable,colback=teal!4,colframe=teal!35!black,title={Before you move on}]
What does dever mean? (\textquotedblleft{}To have to, must\textquotedblright{}.) Say it once more.
\end{tcolorbox}
\section[a necessidade]{a necessidade — a need}

\label{lesson:PT-C82-necessidade}



\begin{quote}
Before the new one: say the Portuguese for to be able to, then the Portuguese for to have to.
\end{quote}

\subsection*{You'll want to know: a necessidade}
\textbf{a necessidade} — \textquotedblleft{}a need\textquotedblright{}.

\begin{cousinweb}
\textbf{necessidade} $\leftarrow$ Latin \emph{necessitas, necessitatis}, \textquotedblleft{}need\textquotedblright{} (EtymDB 2.1) — a learned borrowing from Latin.
\end{cousinweb}

\begin{grammarlens}[title={What you've built}]
One more word for this chapter.
\end{grammarlens}

\subsection*{Your turn}
\begin{itemize}
\raggedright
  \item \emph{Say it:} \emph{a necessidade}
  \item \emph{Say it:} \emph{a necessidade}, once more
  \item \emph{Say it:} say \emph{a necessidade}
  \item \emph{Recall:} say the Portuguese for to be able to, then the Portuguese for to have to, then say \emph{a necessidade} again
\end{itemize}

\begin{tcolorbox}[breakable,colback=teal!4,colframe=teal!35!black,title={Before you move on}]
What does a necessidade mean? (\textquotedblleft{}A need\textquotedblright{}.) Say it once more.
\end{tcolorbox}
\section[o ofício]{o ofício — a trade, a craft}

\label{lesson:PT-C82-oficio}



\begin{quote}
Before the new one: say the Portuguese for to have to, then the Portuguese for a need.
\end{quote}

\subsection*{You'll want to know: o ofício}
\textbf{o ofício} — \textquotedblleft{}a trade, a craft\textquotedblright{}.

\begin{cousinweb}
\textbf{ofício} $\leftarrow$ Latin \emph{officium}, \textquotedblleft{}duty, service\textquotedblright{} (EtymDB 2.1) — a learned borrowing from Latin.
\end{cousinweb}

\begin{grammarlens}[title={What you've built}]
One more word for this chapter.
\end{grammarlens}

\subsection*{Your turn}
\begin{itemize}
\raggedright
  \item \emph{Say it:} \emph{o ofício}
  \item \emph{Say it:} \emph{o ofício}, once more
  \item \emph{Say it:} say \emph{o ofício}
  \item \emph{Recall:} say the Portuguese for to have to, then the Portuguese for a need, then say \emph{o ofício} again
\end{itemize}

\begin{tcolorbox}[breakable,colback=teal!4,colframe=teal!35!black,title={Before you move on}]
What does o ofício mean? (\textquotedblleft{}A trade, a craft\textquotedblright{}.) Say it once more.
\end{tcolorbox}
\section[levar]{levar — to take, to carry}

\label{lesson:PT-C82-levar}



\begin{quote}
Before the new one: say the Portuguese for a need, then the Portuguese for a trade.
\end{quote}

\subsection*{You'll want to know: levar}
\textbf{levar} — \textquotedblleft{}to take, to carry\textquotedblright{}.

\begin{cousinweb}
\textbf{levar} $\leftarrow$ Latin \emph{levare}, \textquotedblleft{}to lift\textquotedblright{} (EtymDB 2.1), from \emph{levis}, \textquotedblleft{}light\textquotedblright{} — inherited.
\end{cousinweb}

\begin{grammarlens}[title={What you've built}]
One more word for this chapter.
\end{grammarlens}

\subsection*{Your turn}
\begin{itemize}
\raggedright
  \item \emph{Say it:} \emph{levar}
  \item \emph{Say it:} \emph{levar}, once more
  \item \emph{Say it:} say \emph{levar}
  \item \emph{Recall:} say the Portuguese for a need, then the Portuguese for a trade, then say \emph{levar} again
\end{itemize}

\begin{tcolorbox}[breakable,colback=teal!4,colframe=teal!35!black,title={Before you move on}]
What does levar mean? (\textquotedblleft{}To take, to carry\textquotedblright{}.) Say it once more.
\end{tcolorbox}
